\subsection{বায়োইনফরমেটিক্স}

\begin{learningbox}
এই অংশ শেষে শিক্ষার্থী পারবে:
\begin{itemize}
\item বায়োইনফরমেটিক্সের ধারণা ব্যাখ্যা করতে।
\item DNA, protein sequence, database ও computational analysis-এর সম্পর্ক বুঝতে।
\item চিকিৎসা, কৃষি ও গবেষণায় বায়োইনফরমেটিক্সের ব্যবহার লিখতে।
\end{itemize}
\end{learningbox}

\subsubsection{বায়োইনফরমেটিক্সের ধারণা}

জীববিজ্ঞানের গবেষণায় DNA sequence, protein structure, gene expression, disease data এবং organism information বিপুল পরিমাণে তৈরি হয়। এই biological data computer, database, algorithm ও statistical method দিয়ে সংরক্ষণ, বিশ্লেষণ ও ব্যাখ্যা করার ক্ষেত্র হলো বায়োইনফরমেটিক্স।

\begin{definitionbox}{সংজ্ঞা}
বায়োইনফরমেটিক্স হলো জীববিজ্ঞান, কম্পিউটার বিজ্ঞান, গণিত ও পরিসংখ্যানের সমন্বিত ক্ষেত্র, যেখানে biological data সংরক্ষণ, অনুসন্ধান, বিশ্লেষণ ও ব্যাখ্যার জন্য computational পদ্ধতি ব্যবহার করা হয়।
\end{definitionbox}

NCBI ও NIH-ভিত্তিক bioinformatics resources-এ sequence database, alignment, genome analysis এবং protein information গুরুত্বপূর্ণ বিষয় হিসেবে বিবেচিত।

\subsubsection{জৈব তথ্য সংরক্ষণ}

Biological data database-এ সংরক্ষণ করা হয়, যাতে researcher দ্রুত data খুঁজে পেতে পারে। DNA sequence database, protein database, gene expression database এবং disease database গবেষণায় ব্যবহৃত হয়।

\begin{keypointbox}{জৈব তথ্যের ধরন}
\begin{itemize}
\item DNA ও RNA sequence।
\item protein sequence ও structure।
\item gene ও chromosome information।
\item organism classification data।
\item disease mutation ও drug response data।
\end{itemize}
\end{keypointbox}

\subsubsection{ডিএনএ বিশ্লেষণ}

DNA sequence analysis করে gene শনাক্ত, mutation খোঁজা, রোগের genetic কারণ বোঝা এবং organism-এর evolutionary relation বিশ্লেষণ করা যায়। Sequence alignment হলো দুটি বা ততোধিক sequence তুলনা করার সাধারণ পদ্ধতি।

\subsubsection{প্রোটিন বিশ্লেষণ}

Protein analysis-এ protein-এর sequence, folding, structure ও function বোঝার চেষ্টা করা হয়। Drug design ও disease research-এ protein structure গুরুত্বপূর্ণ।

\subsubsection{চিকিৎসা ও গবেষণায় ব্যবহার}

\begin{examplebox}{চিকিৎসা গবেষণায় ব্যবহার}
কোনো রোগীর gene mutation জানা থাকলে চিকিৎসক personalized medicine বা targeted therapy নিয়ে ভাবতে পারেন। আবার virus genome analysis করে variant শনাক্ত করা যায়।
\end{examplebox}

\begin{keypointbox}{ব্যবহারক্ষেত্র}
\begin{itemize}
\item disease gene শনাক্তকরণ।
\item vaccine ও drug design research।
\item pathogen বা virus variant analysis।
\item forensic DNA analysis।
\item evolutionary biology research।
\end{itemize}
\end{keypointbox}

\subsubsection{কৃষি গবেষণায় ব্যবহার}

ফসলের disease resistance, drought tolerance, yield improvement এবং genetic diversity বোঝার জন্য bioinformatics ব্যবহৃত হয়। উদ্ভিদের genome data বিশ্লেষণ করে উন্নত জাত উদ্ভাবনে সহায়তা করা যায়।

\begin{mistakebox}{সাধারণ ভুল}
Bioinformatics শুধু biology নয়, আবার শুধু computer science-ও নয়। এটি biological problem সমাধানে computational method ব্যবহারের ক্ষেত্র।
\end{mistakebox}

\begin{boardbox}{বোর্ড ফোকাস}
বায়োইনফরমেটিক্সের প্রশ্নে “biological data + computer analysis” এই সম্পর্কটি পরিষ্কার করো। DNA, protein, database, drug design ও কৃষি গবেষণা—এই শব্দগুলো exam-friendly।
\end{boardbox}

\begin{practicebox}
\textbf{MCQ}
\begin{enumerate}
\item DNA sequence বিশ্লেষণে computer ব্যবহার কোন ক্ষেত্রের উদাহরণ?\\
ক. Bioinformatics \quad খ. Desktop publishing \quad গ. Animation only \quad ঘ. Printing
\item Sequence alignment কী কাজে ব্যবহৃত হয়?\\
ক. sequence তুলনা \quad খ. monitor cleaning \quad গ. file delete \quad ঘ. font install
\end{enumerate}

\textbf{সংক্ষিপ্ত প্রশ্ন}
\begin{enumerate}
\item বায়োইনফরমেটিক্স কী?
\item কৃষি গবেষণায় বায়োইনফরমেটিক্সের দুটি ব্যবহার লেখ।
\end{enumerate}
\end{practicebox}

